\documentclass[report]{itvest}

% Document-specific metadata
\title{Portfolio 3: Big Data \& Data Science Technologies}
\reportsubtitle{Case: Short Anonymised Case Name}
\author{Student Name} % group: \author{Name 1 \and Name 2 \and Name 3}
\date{\today}

% Portfolios use sections only, numbered 1, 2, ...
\renewcommand{\thesection}{\arabic{section}}

\begin{document}

\maketitle

\begin{checkit}
  \textbf{Assessment criteria} (delete this box before submitting):
  \begin{itemize}
    \item The volume estimate is calculated and drives the processing and
      storage choices.
    \item Governance rules are concrete and consistent with the P2
      data-placement table.
    \item Drift and quality signals are defined, measurable and tied to an
      action.
  \end{itemize}
\end{checkit}

\section{Data Flow}

Show how data flows from source to storage to model, with the tier
boundaries from P2.

\begin{figure}[h]
  \centering
  % Replace the box with: \includegraphics[width=\textwidth]{data-flow}
  \fbox{\parbox[c][7cm][c]{0.9\textwidth}{\centering
    Replace with the data-flow diagram, tier boundaries marked.}}
  \caption{Data flow.}
  \label{fig:data-flow}
\end{figure}

\section{Volume and Rate}

Estimate what data the system produces and consumes, at what volume and
rate. Show the calculation, e.g.
\[
  \text{volume per day} =
  \text{sources} \times \text{records per second} \times
  \text{bytes per record} \times 86\,400 .
\]

\begin{table}[h]
  \centering
  \begin{tabularx}{\textwidth}{Xlll}
    \toprule
    Data source & Rate & Size per record & Volume per day \\
    \midrule
    \ldots & \ldots & \ldots & \ldots \\
    \ldots & \ldots & \ldots & \ldots \\
    \bottomrule
  \end{tabularx}
  \caption{Volume and rate estimate.}
  \label{tab:volume}
\end{table}

\section{Processing and Storage}

Choose batch, stream or both, and where historical and event data are
stored (topic retention, data lake, warehouse, feature
store)~\cite{huyen_designing_2022}. Use the estimate in
Table~\ref{tab:volume} as an argument.

\subsection{ADR-4: Processing}

\begin{description}
  \item[Status:] Proposed / accepted.
  \item[Context:] The volume, rate and latency that drive the choice.
  \item[Alternatives:] Batch, stream, or both.
  \item[Decision:] The chosen alternative.
  \item[Consequences:] The trade-off that decided it, and what gets harder.
\end{description}

\subsection{ADR-5: Storage}

\begin{description}
  \item[Status:] \ldots
  \item[Context:] \ldots
  \item[Alternatives:] \ldots
  \item[Decision:] \ldots
  \item[Consequences:] \ldots
\end{description}

\section{Data Governance}

State who owns which data, how long it is kept, and which regulations
apply (e.g.\ GDPR, NIS2, EU AI Act), and what they require of the design.
The locations must match the P2 data-placement table.

\begin{table}[h]
  \centering
  \begin{tabularx}{\textwidth}{XlllX}
    \toprule
    Data item & Owner & Retention & Location & Applicable rule \\
    \midrule
    \ldots & \ldots & \ldots & \ldots & \ldots \\
    \ldots & \ldots & \ldots & \ldots & \ldots \\
    \bottomrule
  \end{tabularx}
  \caption{Data governance.}
  \label{tab:governance}
\end{table}

\section{Observability}

\subsection{Data Quality}

Describe the checks on incoming data (e.g.\ schema, ranges, missing
values) and where they run.

\subsection{Drift and Proxy Metrics}

Describe how drift is detected while labels arrive late, and which proxy
metrics stand in for model quality until they do.

\subsection{Service and Model Quality}

Separate service SLOs from model-quality metrics.

\begin{table}[h]
  \centering
  \begin{tabularx}{\textwidth}{XllX}
    \toprule
    Metric & Kind & Target & Measured where \\
    \midrule
    \ldots & Service SLO & \ldots & \ldots \\
    \ldots & Model quality & \ldots & \ldots \\
    \ldots & Data quality / drift & \ldots & \ldots \\
    \bottomrule
  \end{tabularx}
  \caption{Observability metrics.}
  \label{tab:metrics}
\end{table}

\subsection{Feedback Loop}

Tie each signal to an action: retrain, roll back or reconfigure.

\begin{table}[h]
  \centering
  \begin{tabularx}{\textwidth}{XXX}
    \toprule
    Signal & Threshold & Action \\
    \midrule
    \ldots & \ldots & \ldots \\
    \bottomrule
  \end{tabularx}
  \caption{Signals and actions.}
  \label{tab:feedback}
\end{table}

\section{Evidence}

State whether the evidence comes from your own case or from the reference
project, and show the running pipeline. Dashboards and drift evidence from
week 6 may be added in the final report.

\section{AI Usage}

Briefly describe your use of AI tools on this assignment:

\begin{itemize}
  \item \emph{Code generation:} \ldots
  \item \emph{Grammar and writing:} \ldots
  \item \emph{Clarification of technical aspects:} \ldots
\end{itemize}

\bibliographystyle{IEEEtran}
\bibliography{bib}

\end{document}
